
\section{New Material}

Consider this hypothetical announcement: 

\begin{equation}
\label{badAnnouncement}
\parbox{10cm}{\textit{Your flight from Manchester to Venice will leave 
from Manchester terminal 3 }}
\end{equation}

There is something badly formed about it in that the second occurrence of the word Manchester is redundant --- it can be announced instead that
\begin{equation}
\label{goodAnnouncement}
\parbox{10cm}{\textit{Your flight from Manchester to Venice will leave 
from terminal 3}} 
\end{equation}
for it is a tautology that terminal 3 must be the terminal 3 of Manchester airport. 
These announcements are everyday language analogues of data. In relational data theory the analogue of (\ref{badAnnouncement}) is said to fail the 3rd normal form criteria whereas (\ref{goodAnnouncement}) is said to be in normal form. Here in this book instead of talking of normal forms we shall prefer to speak of goodness principles. These, analogues of (\ref{badAnnouncement}) fail whilst analogues of (\ref{goodAnnouncement}) pass. This is a different approach to the goodness of data and one that has far more inherent sense than the seemingly ad hoc normal form criteria of relational data theory.

That there is a flight from Manchester to Venice leaving terminal 3 is particular knowledge and is the stuff of data. That, generally, 
\begin{equation}
\label{leavingTerminalRule}
 \parbox{10cm}{\textit{a flight leaving an airport leaves from a terminal of that same airport }}
 \end{equation}
 is for all practical purposes, universal knowledge. 
 It is a general rule. 
 The imperfection in the particular statement (\ref{badAnnouncement}) is that the way  formulated does not take advantage of the universal 
 (\ref{leavingTerminalRule}).

\begin{worktt}
Courtesy of ChatGPT:

* Treat ordinary declarative language as a primary form of data, not merely an analogy for data.
* Introduce data via particular knowledge rather than via database records.
* Distinguish particular knowledge from universal knowledge.
* Explain redundancy as failure to make proper use of available universal knowledge.
* Present normal forms as a special case of more general "goodness principles".
* Explain database representations as one representation of declarative information among others.
* Use the universal/particular distinction to explain why there are two broad kinds of type (entity types and attribute/domain types).
* Ground terminology in established mathematical and philosophical language where possible, to avoid reinvention.

One thing that strikes me in retrospect is that your Manchester–Venice example may be more important than it first appears. It is not merely illustrating redundancy. It is illustrating a phenomenon that exists *before databases*. The reader can recognise the defect without knowing anything about relations, keys, or normal forms.

That is quite powerful because it reverses the usual pedagogical direction. Instead of saying:

> Here is a database rule; now let me illustrate it with an English sentence.

you are effectively saying:

> Here is a defect in the expression of knowledge that everyone can recognise; normalisation theory is one way of dealing with the same phenomenon in databases.

That is a much deeper claim.

And it ties rather nicely to something you've often emphasised: conceptual modelling should not be subordinate to implementation technology. If the phenomenon exists already in ordinary language, then the database theory is describing a special case of a more general conceptual phenomenon.

So I would definitely make a note of that before diving back into range categories. It feels less like a presentational trick and more like a genuine shift in foundations. The theoretical paper and the book may ultimately be approaching the same mountain from opposite sides:

* the paper starts with categorical structure and works toward data;
* the book may start with ordinary declarative knowledge and work toward data.

Those are surprisingly complementary viewpoints.

\end{worktt}